\hypertarget{hackstreak-spec-v2-hackstreak-v2-streaks--editable-roadmaps-packs--optional-auto-verification}{%
\subsection{hackstreak v2: streaks + editable roadmaps ("packs") +
optional
auto-verification}\label{hackstreak-spec-v2-hackstreak-v2-streaks--editable-roadmaps-packs--optional-auto-verification}}

\textbf{Precedence:} this file supersedes \texttt{hackstreak-spec.md}
(v1) wherever they conflict. Keep v1 sections 6 (rules), 8 (streak
algorithm), 10 (security) and the streak tests in 11 exactly as written.
Replaced by this file: v1 sections 4, 5, 7, 9 and 12.

\begin{center}\rule{0.5\linewidth}{0.5pt}\end{center}

\hypertarget{hackstreak-spec-v2-1-the-design-in-one-paragraph}{%
\subsubsection{1. The design in one
paragraph}\label{hackstreak-spec-v2-1-the-design-in-one-paragraph}}

The app has three layers. The \textbf{engine} (code) knows streaks,
logging and "what is my next step", and knows nothing about any subject.
\textbf{Packs} (plain TOML text files the owner can edit) hold the
roadmap: nodes with prerequisites, size, a real resource link and a
mandatory "done when" sentence. \textbf{Verifiers} (small plugins) can
check a node automatically (a file exists, a command succeeds,
Codeforces shows a new accepted problem). The map is therefore DATA the
owner controls, not behaviour hardcoded in the program, and auto-fetch
is an optional bonus on top of honest manual marking.

\hypertarget{hackstreak-spec-v2-2-principles-added-in-v2}{%
\subsubsection{2. Principles added in
v2}\label{hackstreak-spec-v2-2-principles-added-in-v2}}

\begin{enumerate}
\def\labelenumi{\arabic{enumi}.}
\tightlist
\item
  \textbf{Map = data.} Nothing about Bandit, Juice Shop or any platform
  appears in engine code. Ever.
\item
  \textbf{Packs are read-only inputs; progress is separate.} Pack files
  are never modified by the app. Progress lives in SQLite. Updating a
  pack never loses progress (nodes are matched by \texttt{id}).
\item
  \textbf{Every node has a \texttt{done\_when}.} An ambiguous "did I do
  it?" is a pack-authoring defect and \texttt{lint} rejects it.
\item
  \textbf{Streak day and node completion are different things.} Any real
  work counts for the streak (partial progress is fine). A node
  completes only when its \texttt{done\_when} is met. This removes the
  "I did something but can\textquotesingle t tick anything" feeling.
\item
  \textbf{Verifiers propose, humans confirm,} except hard local checks,
  which may auto-complete if the node allows it.
\item
  \textbf{Generated maps are drafts.} A pack drafted by an AI is
  labelled \texttt{source = "ai-draft"} and
  \texttt{reviewed = false} and every node it shows is marked
  \texttt{{[}draft{]}} until the owner runs \texttt{hack review}. AI
  may order and link; it must not invent lessons or URLs.
\end{enumerate}

\hypertarget{hackstreak-spec-v2-3-pack-format-toml-readable-with-pythons-standard-tomllib}{%
\subsubsection{\texorpdfstring{3. Pack format (TOML, readable with
Python\textquotesingle s standard
\texttt{tomllib})}{3. Pack format (TOML, readable with Python\textquotesingle s standard tomllib)}}\label{hackstreak-spec-v2-3-pack-format-toml-readable-with-pythons-standard-tomllib}}

\begin{verbatim}
[pack]
id = "home-hack-lab"          # unique, kebab-case
title = "Home hack lab"
version = 1
source = "ai-draft"           # human | ai-draft | imported
reviewed = false

[vars]                        # user-editable values, substituted as {name} inside strings
pi_user = "pi"
pi_ip = "192.168.1.50"

# Levels group nodes (like Duolingo units). Order in the file = order of the main path.
[[level]]
id = "setup"
title = "Build the lab"
lane = "main"                 # main = gated path (default) | side = always-open quick wins
gate = true                   # main lanes only: level opens when the previous main level is fully done/skipped

[[level]]
id = "warmup"
title = "Daily warm-up"
lane = "side"                 # side lanes never block and never wait; they supply the "quick win" for busy days

[[node]]
level = "setup"
id = "pi-reachable"           # unique within the pack, stable forever (progress is keyed on it)
title = "Pi is reachable from the laptop"
size = "M"                    # S <= 30 min, M = 1-2 h, L = 3 h+
needs = []                    # prerequisite node ids (must form a DAG)
done_when = "ping to the Pi works from the laptop in both directions"
resource = "https://www.raspberrypi.com/documentation/computers/remote-access.html"
verify = { type = "command", cmd = "ping -c 1 -W 2 {pi_ip}" }   # optional; default is manual
auto_complete = false         # true = a passing verifier completes the node without asking

# Series: one definition that expands into many chained nodes (generic, works for any numbered material).
[[node]]
id = "bandit"
title = "OverTheWire Bandit, level {n}"
series = { from = 13, to = 33 }
size = "S"
needs = []
done_when = "I logged in to level {next} and can explain the command that got me there"
resource = "https://overthewire.org/wargames/bandit/"
\end{verbatim}

Rules enforced by
\path|hack pack lint <file>|:

\begin{itemize}
\tightlist
\item
  \texttt{id}s unique; \texttt{needs} reference existing ids; graph is
  acyclic.
\item
  every node has a \texttt{level} that exists; levels have \texttt{id},
  \texttt{title}, \texttt{lane} (\texttt{main} default, or
  \texttt{side}); at least one \texttt{main} level exists.
\item
  a \texttt{series} node lives in exactly one level (use two series
  nodes if the material spans levels).
\item
  every node has \texttt{title}, \texttt{size} (S/M/L),
  \texttt{done\_when} (min 15 characters) and \texttt{resource} (an
  http(s) URL).
\item
  series expand to ids \texttt{bandit-13} ... \texttt{bandit-33}; each
  depends on the previous one; the first inherits the
  series\textquotesingle{} \texttt{needs}. \texttt{\{n\}} = current
  number, \texttt{\{next\}} = n + 1.
\item
  unknown \texttt{verify.type} is an error. A pack containing
  \texttt{command} verifiers prints the commands and requires
  \texttt{hack pack trust \textless{}id\textgreater{}} first (see
  section 7).
\item
  \path|hack pack check-links <file>| sends
  a HEAD/GET to every \texttt{resource} and reports dead links. Run it
  on every AI-drafted pack.
\end{itemize}

\hypertarget{hackstreak-spec-v2-4-storage-additions-sqlite}{%
\subsubsection{4. Storage additions
(SQLite)}\label{hackstreak-spec-v2-4-storage-additions-sqlite}}

\begin{verbatim}
CREATE TABLE IF NOT EXISTS packs (
  id TEXT PRIMARY KEY, path TEXT NOT NULL, version INTEGER, reviewed INTEGER DEFAULT 0, trusted INTEGER DEFAULT 0
);
CREATE TABLE IF NOT EXISTS node_state (
  pack_id TEXT NOT NULL, node_id TEXT NOT NULL,
  status TEXT NOT NULL DEFAULT 'todo',      -- todo | done | skipped
  done_at TEXT, evidence TEXT, how TEXT,    -- how: manual | verifier:<type> | forced
  PRIMARY KEY (pack_id, node_id)
);
-- entries (from v1) gain three columns:
ALTER TABLE entries ADD COLUMN node_id TEXT;       -- "pack:node", nullable
ALTER TABLE entries ADD COLUMN source TEXT DEFAULT 'manual';   -- manual | verifier:<type>
ALTER TABLE entries ADD COLUMN complete INTEGER DEFAULT 0;
-- missions (from v1) gain: pack_id TEXT  (a mission = "work through this pack"); still only ONE active mission.
\end{verbatim}

\hypertarget{hackstreak-spec-v2-5-choosing-the-next-step-the-zero-decisions-rule}{%
\subsubsection{5. Choosing the next step (the "zero decisions"
rule)}\label{hackstreak-spec-v2-5-choosing-the-next-step-the-zero-decisions-rule}}

\begin{verbatim}
SIZE_MIN = {"S": 30, "M": 120, "L": 240}

def open_levels(levels, nodes, finished):
    opened, prev_complete = set(), True
    for lv in levels:                                  # pack order
        if lv.lane == "side":
            opened.add(lv.id)                          # side lanes are always open
            continue
        if prev_complete or not lv.gate:
            opened.add(lv.id)
        prev_complete = all(n.id in finished for n in nodes if n.level == lv.id)
    return opened

def next_nodes(nodes, levels, state, max_minutes=None):
    finished = {nid for nid, st in state.items() if st in ("done", "skipped")}
    opened = open_levels(levels, nodes, finished)
    out = []
    for n in nodes:                                    # pack order = tie-breaker
        if n.id in finished:
            continue
        if n.level not in opened:                      # level gate
            continue
        if not all(r in finished for r in n.needs):    # locked by prerequisites
            continue
        # NOTE: prev_complete in open_levels() checks only the milestone node(s) of a level when any exist,
        # not every node in it -- gate = "passed the level's real test", not "touched every item".
        if max_minutes is not None and SIZE_MIN[n.size] > max_minutes:
            continue
        out.append(n)
    return out
\end{verbatim}

\begin{itemize}
\tightlist
\item
  \texttt{hack} shows the first candidate from a \texttt{main} lane as
  \textbf{Next}, and the first candidate from a \texttt{side} lane as
  \textbf{Quick win} (so a busy day always has a 10-minute option). No
  full list by default; \texttt{hack more} shows two alternatives.
\item
  Level display:
  \path|Level 2/3  Recon & attack  {[}####-----------{]} 3/8 steps  5 active days|.
  "Active days" = distinct logical days with an entry whose node belongs
  to that level. This is the per-level, day-wise view; the global daily
  streak is unchanged.
\item
  \texttt{hack levels} lists every level with progress and active days.
  \texttt{hack unlock \textless{}level\textgreater{}} opens a gated
  level early (logged), so a gate can never trap the owner.
  \texttt{hack skip \textless{}node\textgreater{}} is
  \texttt{mark ... skipped}.
\item
  With \texttt{hack \textless{}minutes\textgreater{}}: show the
  main-lane step if it fits, otherwise the side-lane quick win,
  otherwise suggest \texttt{hack split}.
\item
  \texttt{hack 20} = "I have 20 minutes": only nodes that fit. If none
  fit it says so and suggests \texttt{hack split}.
\item
  On an academic/deadline day the owner types \texttt{hack 10} and gets
  the smallest valid step. The streak survives.
\end{itemize}

\hypertarget{hackstreak-spec-v2-6-entry-semantics-partial-vs-complete}{%
\subsubsection{6. Entry semantics (partial vs
complete)}\label{hackstreak-spec-v2-6-entry-semantics-partial-vs-complete}}

\begin{itemize}
\tightlist
\item
  \path|hack done "learned..." "next step..."| logs an entry
  (counts for the streak) against the current node. The node stays open.
\item
  \path|hack done ... ---complete| also tries to mark the node
  done: it runs the node\textquotesingle s verifier first.

  \begin{itemize}
  \tightlist
  \item
    verifier passes (or node is manual): node becomes \texttt{done},
    \texttt{how} = manual or verifier name.
  \item
    verifier fails: print its reason, do NOT complete.
    \texttt{-\/-force "reason"} overrides and records
    \texttt{how = forced}.
  \end{itemize}
\item
  \texttt{hack check} runs the current node\textquotesingle s verifier
  on its own and, if it passes, offers to complete it.
\item
  \path!hack mark <node> done|skipped|todo {[}---note ...{]}!
  is the honest manual route. It is always allowed.
\item
  \path|hack backfill <node> {[}<node>...{]}|
  marks nodes done with \texttt{how = manual},
  \path|evidence = "completed before hackstreak"|, without
  creating streak entries. (The owner uses this for Bandit 0-12.)
\end{itemize}

\hypertarget{hackstreak-spec-v2-7-verifiers-plugins}{%
\subsubsection{7. Verifiers (plugins)}\label{hackstreak-spec-v2-7-verifiers-plugins}}

Interface, one small file per verifier in
\path|hackstreak/verifiers/|:

\begin{verbatim}
class VerifyResult(NamedTuple):
    ok: bool
    summary: str            # shown to the user and stored as evidence, e.g. "nmap output 1.2 KB, modified today"
    evidence_time: datetime | None   # when the real-world thing happened (None = unknown)

def verify(spec: dict, vars: dict) -> VerifyResult: ...
\end{verbatim}

Built-ins by phase:

\begin{longtable}[]{@{}P{.24\textwidth}P{.35\textwidth}P{.33\textwidth}@{}}
\toprule\noalign{}
Type & Checks & Phase \\
\midrule\noalign{}
\endhead
\bottomrule\noalign{}
\endlastfoot
\texttt{manual} & nothing; human confirms & 1 \\
\texttt{file} & path exists, optional \texttt{min\_bytes}, optional
\texttt{contains} regex & 1 \\
\texttt{command} & runs a shell command, exit code 0 = pass, 10 s
timeout, output truncated & 1 \\
\texttt{http} & GET url, expected status & 2 \\
\texttt{git} & repo has a commit matching a regex since a date & 2 \\
\texttt{codeforces} & public API: accepted problems since node start
\textgreater= N & 3 \\
\texttt{praxis} & reads the local Praxis API (\texttt{/api/overview},
\texttt{/api/submissions}) for solved-count deltas & 3 \\
\texttt{github} & public API: commits/pushes since a date & 3 \\
\end{longtable}

\textbf{Trust model.} \texttt{command} verifiers execute shell commands
from a pack file. A pack from someone else (or from an AI) is code.
Therefore: \texttt{hack pack install} prints every command, nothing
runs until \texttt{hack pack trust \textless{}id\textgreater{}} is
typed, and editing a pack resets trust. Commands run with the
owner\textquotesingle s permissions, with no \texttt{sudo}, never
through a shell string built from user input (use \texttt{shlex.split},
\texttt{shell=False}).

\textbf{Auto evidence counts.} When a verifier passes with a known
\texttt{evidence\_time}, the app writes an entry for that day with
\path|source = "verifier:<type>"| and an
auto-generated \texttt{learned} line (the verifier summary). Proof is
the verifier, so the 20-character rule applies only to manual entries.
Setting \path|auto_counts_for_streak| (default true) can turn this
off.

\texttt{hack sync} runs the verifiers of all UNLOCKED nodes, lists what
looks done, and either completes (\texttt{auto\_complete = true}) or
asks \texttt{hack accept}.

\hypertarget{hackstreak-spec-v2-8-what-can-honestly-be-auto-fetched-verify-these-before-building-apis-change}{%
\subsubsection{8. What can honestly be auto-fetched (verify these before
building, APIs
change)}\label{hackstreak-spec-v2-8-what-can-honestly-be-auto-fetched-verify-these-before-building-apis-change}}

\begin{longtable}[]{@{}P{.22\textwidth}P{.18\textwidth}P{.32\textwidth}P{.18\textwidth}@{}}
\toprule\noalign{}
Source & Auto-detect? & How & My confidence \\
\midrule\noalign{}
\endhead
\bottomrule\noalign{}
\endlastfoot
Anything on the owner\textquotesingle s machines (files, scans,
containers, ssh works, git commits) & Yes & \texttt{file} /
\texttt{command} / \texttt{git} / \texttt{http} & High, fully under the
owner\textquotesingle s control \\
Codeforces & Yes & public API & High \\
GitHub & Yes & public API & High \\
Praxis (own app) & Yes & already syncs CF/LC/AtCoder/GfG, reuse its API
& High \\
LeetCode, GfG & Partly & unofficial endpoints, may break & Medium \\
TryHackMe, HackTheBox & Unclear & no stable official progress API that I
am sure of; check docs first & Low \\
OverTheWire (Bandit etc.), picoCTF, books, videos & No & no
account-progress API & Use manual + a concrete \texttt{done\_when} \\
\end{longtable}

\textbf{Artifact trick for things that cannot be fetched:} make the
proof something that exists on the owner\textquotesingle s disk (a short
writeup file, a scan, a script) and verify THAT with the \texttt{file}
verifier. It also builds the habit of explaining things. \textbf{Never
store or log wargame passwords or flags} (v1 secrets guard stays).

\hypertarget{hackstreak-spec-v2-8b-pack-scale-project-packs-vs-proficiency-packs}{%
\subsubsection{8b. Pack scale: "project" packs vs "proficiency"
packs}\label{hackstreak-spec-v2-8b-pack-scale-project-packs-vs-proficiency-packs}}

A pack that finishes in a week solves nothing; the owner is back to
"what now" within days. Two kinds of pack, named explicitly so neither
gets mistaken for the other:

\begin{itemize}
\tightlist
\item
  \textbf{Project pack} (\path|pack.kind = "project"|): a concrete
  build with a finish line, 3-6 levels, finishes in days-to-weeks.
  Example: the home hack lab itself.
\item
  \textbf{Proficiency pack} (\path|pack.kind = "proficiency"|):
  meant to make the owner genuinely good at a subject (networking, OS
  internals, a language, DSA topic, web/mobile dev). \textbf{8-15 main
  levels}, each a real sub-skill, spanning weeks to a few months at the
  owner\textquotesingle s real pace. This is the default for "I want to
  learn X" requests.
\end{itemize}

Every level in a proficiency pack ends with one \textbf{milestone node}
(\texttt{milestone = true}, size \texttt{M} or \texttt{L}): a harder,
integrative task that cannot be passed by rote repetition of the
level\textquotesingle s small nodes, e.g. "build a tiny TCP chat
client+server from scratch, no copied code, and explain every socket
call" rather than "read about sockets." A level\textquotesingle s gate
(section 5) checks on the milestone node specifically, not merely on
every node being done, so "proficient in this level" has a real test
behind it, the same way a Duolingo unit ends with a harder review
lesson. \texttt{lint} requires at least one \texttt{milestone = true}
node per main level.

\hypertarget{hackstreak-spec-v2-9-pack-authoring-how-maps-get-made-without-losing-control}{%
\subsubsection{9. Pack authoring (how maps get made without losing
control)}\label{hackstreak-spec-v2-9-pack-authoring-how-maps-get-made-without-losing-control}}

Three sources, all ending in the same lint -\textgreater{} review
-\textgreater{} install pipeline:

\begin{enumerate}
\def\labelenumi{\arabic{enumi}.}
\tightlist
\item
  \textbf{By hand:} \path|hack pack new "Topic"| scaffolds a TOML
  file with one example node.
\item
  \textbf{AI draft:} the owner (or Claude, or a local model) produces
  TOML from the prompt below, then: \texttt{hack pack lint}
  -\textgreater{} \path|hack pack check-links| -\textgreater{} read
  it -\textgreater{}
  \texttt{hack review \textless{}pack\textgreater{}} (flips
  \texttt{reviewed}) -\textgreater{} \texttt{hack pack install}.
\item
  \textbf{Imported:} converted from a public roadmap by hand. Check that
  source\textquotesingle s licence before copying its content.
\end{enumerate}

Prompt for generating a pack draft (paste into any LLM):

\begin{quote}
Output ONLY a TOML file in the hackstreak pack format (sections
\texttt{{[}pack{]}}, \texttt{{[}vars{]}}, \texttt{{[}{[}node{]}{]}}).
Topic: . My level: . Hardware: . Weekly time: . Rules: every node needs
id, title, size (S\textless=30min, M=1-2h, L=3h+), needs, done\_when
(one concrete checkable sentence, ideally an artifact on my disk), and
resource (a real, well-known URL; if you are not sure a URL exists,
write \texttt{resource = "TODO"} instead of guessing). Do not write
lesson content. Prefer free official or community resources. Use
\texttt{series} for numbered material. If this is a PROFICIENCY request
(learn a subject, not build one project), use
\path|pack.kind = "proficiency"| and write 8-15
\texttt{{[}{[}level{]}{]}} blocks on the \texttt{main} lane (like
Duolingo units), each 4-8 nodes covering one real sub-skill, and end
each level with one \texttt{milestone = true} node that is a harder
integrative task (build/explain something from scratch), not a repeat of
the level\textquotesingle s small steps. If this is a PROJECT request
(one concrete build with a finish line), use
\path|pack.kind = "project"| and keep it to 3-6 levels total. Either
way add at most one \texttt{side} lane level of S-sized quick wins for
busy days. Every node needs a \texttt{level}. Order so that the first
node can be started tonight. Prefer done\_when sentences that produce an
artifact on my disk. Set \texttt{source = "ai-draft"} and
\texttt{reviewed = false}.
\end{quote}

\hypertarget{hackstreak-spec-v2-9b-workflow-ask-claude-for-each-new-unit}{%
\paragraph{9b. Workflow: "ask Claude for each new
unit"}\label{hackstreak-spec-v2-9b-workflow-ask-claude-for-each-new-unit}}

The engine ships with one starter pack. Every later unit is a new pack
drafted on request:

\begin{enumerate}
\def\labelenumi{\arabic{enumi}.}
\tightlist
\item
  Fill in \path|new-unit-request-template.md| (topic, current level,
  time, hardware, what is already done) and paste it into a chat.
\item
  The model replies with ONLY a TOML pack. Save it to
  \texttt{packs/\textless{}id\textgreater{}.toml}.
\item
  \texttt{hack pack lint} -\textgreater{}
  \path|hack pack check-links| -\textgreater{} read every
  \texttt{done\_when} -\textgreater{}
  \texttt{hack review \textless{}id\textgreater{}} -\textgreater{}
  \path|hack pack install <id>|.
\item
  \texttt{hack backfill ...} for anything already done, then
  \path|hack mission start <id>| (only one
  active mission at a time). The model has no memory between chats, so
  the template carries everything it needs; do not rely on it
  remembering the format.
\end{enumerate}

\hypertarget{hackstreak-spec-v2-10-cli-summary-v1-commands-plus}{%
\subsubsection{10. CLI summary (v1 commands
plus)}\label{hackstreak-spec-v2-10-cli-summary-v1-commands-plus}}

\begin{verbatim}
hack                         next node, mission, streak, today status (max 6 lines)
hack 20                      next node that fits in 20 minutes
hack more                    two alternatives to the current next node
hack done "..." "..." [--complete] [--force "why"] [--proof X]
hack check                   run current node's verifier
hack sync                    run verifiers on all unlocked nodes, propose completions
hack accept                  confirm proposed completions
hack mark <node> done|skipped|todo [--note ...]
hack backfill <node>...      mark old progress without creating streak entries
hack pack new|lint|check-links|install|trust|list|diff|review
hack mission start <pack>    make a pack the active mission (one active at a time)
hack levels                  every level: progress bar, steps done/total, active days
hack unlock <level>          open a gated level early (always allowed, always logged)
hack skip <node>             shorthand for: mark <node> skipped
hack split <node> "a" "b"    (Phase 2) replace a node with owner-defined smaller steps, stored in the DB, pack untouched
hack park / parked / stats / heatmap / export / remind-check     (unchanged from v1)
hack config nudge-time HH:MM | nudge-channel telegram|desktop|both | nudge-test
hack install-nudge            set up the local (Ubuntu) desktop-notification timer
\end{verbatim}

\hypertarget{hackstreak-spec-v2-11-phases-revised}{%
\subsubsection{11. Phases (revised)}\label{hackstreak-spec-v2-11-phases-revised}}

\begin{itemize}
\tightlist
\item
  \textbf{Phase 1 (one sitting, stdlib only):} v1 Phase 1 + packs
  (\texttt{tomllib}), \texttt{lint}, \texttt{install}, \texttt{trust},
  \texttt{mission start}, next-step picking with
  \texttt{hack \textless{}minutes\textgreater{}}, \texttt{mark},
  \texttt{backfill}, \texttt{done -\/-complete}, verifiers
  \texttt{manual} / \texttt{file} / \texttt{command}, milestone-gated
  levels.
\item
  \textbf{Phase 1b (nudging -\/- build right after Phase 1, same week):}
  see section 11b. This exists now, not later, because a system that
  never pings the owner relies on willpower, which is the exact thing
  that has not worked so far.
\item
  \textbf{Phase 2:} \texttt{http} and \texttt{git} verifiers,
  \texttt{sync}/\texttt{accept}, \texttt{split}, \texttt{check-links}.
\item
  \textbf{Phase 3:} Codeforces / Praxis / GitHub verifiers, dashboard,
  shared log across devices.
\item
  \textbf{Phase 4 (deliberately later):} pack sharing, optional
  LeetCode/TryHackMe verifiers if their APIs prove stable.
\end{itemize}

\hypertarget{hackstreak-spec-v2-11b-nudging-phase-1b}{%
\subsubsection{11b. Nudging (Phase 1b)}\label{hackstreak-spec-v2-11b-nudging-phase-1b}}

\textbf{Problem specific to this owner:} the laptop dual-boots Windows
11 and Ubuntu 24.04. A cron job on Ubuntu goes silent on every day spent
in Windows. The Raspberry Pi 5 is the one machine that is always on
regardless of what the laptop is doing, so the Pi is the primary nudge
source. The laptop\textquotesingle s own notifications are a secondary,
bonus channel.

\textbf{Primary: Telegram from the Pi (reliable across both OSes).}

\begin{itemize}
\tightlist
\item
  \texttt{hack sync-to-pi} (or a \texttt{cron}/\texttt{rsync}
  one-liner) mirrors
  \path|~/.hackstreak/hackstreak.db| (or just the
  \texttt{entries}/\texttt{node\_state} tables) to the Pi once a day,
  OR, simpler for Phase 1b: the owner runs \texttt{hack} ad-hoc on the
  Pi itself over SSH, OR the DB lives on the Pi from the start and the
  laptop points \texttt{HACKSTREAK\_HOME} at it over the already-working
  SSH/LAN link from the hack-lab mission. Pick whichever is already true
  by the time this is built; do not add a new sync mechanism just for
  this.
\item
  On the Pi, a systemd timer (or cron) runs at a configurable time
  (\path{settings.nudge_time}, default \texttt{21:30}), executing
  \texttt{hack remind-check}. Exit 1 (not logged today) triggers
  \path|hackstreak/nudge_telegram.py|, which POSTs to
  \path|https://api.telegram.org/bot<token>/sendMessage|
  using \texttt{TELEGRAM\_BOT\_TOKEN}/\texttt{TELEGRAM\_CHAT\_ID} from
  environment variables (never stored in the DB or committed to git;
  \texttt{.env} stays in \texttt{.gitignore}).
\item
  Message includes the streak, "at risk" framing, and the actual next
  step text, e.g.
  \path|"Streak 9 (at risk). Chess done? Next: bandit-14 or recon-scan (20 min)."|
\item
  A second, later check (e.g. 23:30) fires only if still unlogged,
  mirroring Duolingo\textquotesingle s second reminder. No third nudge
  in one day: more than two starts to feel naggy rather than helpful,
  and the owner already named notification fatigue as a risk ("sometimes
  im not able to do this" referred to overload, not forgetting).
\end{itemize}

\textbf{Secondary: local desktop notification on Ubuntu (bonus, same
trigger).}

\begin{itemize}
\tightlist
\item
  \texttt{hack install-nudge} installs a user-level systemd timer
  calling \texttt{hack remind-check}; on failure it shells out to
  \path|notify-send "hackstreak" "<message>"|
  (available by default on Ubuntu/GNOME). This only fires on days the
  owner is booted into Ubuntu at that hour, which is why it is
  secondary, not the thing relied on.
\item
  No equivalent is built for Windows in Phase 1b. If needed later,
  Windows Toast via \texttt{win10toast} or a scheduled task is a Phase
  2+ add-on, not required while the Pi/Telegram path exists.
\end{itemize}

\textbf{Config:} \path|hack config nudge-time 21:30|,
\path!hack config nudge-channel telegram|desktop|both!,
\path|hack config nudge-test| (sends a test message immediately, so
the owner can confirm it works before trusting it).

\textbf{Tests:} \texttt{remind-check} exit code matches whether today is
logged; the nudge script does not fire a second message on the same day
after the owner logs between the two check times; a missing Telegram
token degrades to "desktop only, warn once" rather than crashing
silently.

\hypertarget{hackstreak-spec-v2-12-added-acceptance-tests}{%
\subsubsection{12. Added acceptance
tests}\label{hackstreak-spec-v2-12-added-acceptance-tests}}

\begin{itemize}
\tightlist
\item
  \texttt{lint} rejects: duplicate id, missing \texttt{done\_when},
  \texttt{needs} pointing nowhere, a cycle, size outside S/M/L, non-URL
  resource.
\item
  Series \texttt{\{from=13,to=15\}} expands to \texttt{bandit-13..15},
  chained in order, with \texttt{\{n\}}/\texttt{\{next\}} substituted.
\item
  \texttt{next\_nodes} skips locked nodes, skips done/skipped nodes,
  respects \texttt{max\_minutes}, and returns \texttt{{[}{]}} when
  nothing fits.
\item
  Updating a pack file (adding a node, reordering) keeps existing
  \texttt{node\_state} rows intact.
\item
  \texttt{done -\/-complete} with a failing \texttt{file} verifier
  leaves the node \texttt{todo}; with \texttt{-\/-force} it records
  \texttt{how = forced}.
\item
  \texttt{command} verifier: times out after 10 s, never uses
  \texttt{shell=True}, and a pack with command verifiers cannot run
  before \texttt{trust}.
\item
  \texttt{backfill} creates \texttt{node\_state} rows but zero entries
  (streak unchanged).
\item
  Verifier-sourced entries count for the streak on the day of
  \texttt{evidence\_time}, not the day of \texttt{sync}.
\item
  Levels: nodes in main level 2 are hidden until every node of main
  level 1 is done/skipped; \texttt{unlock} opens it early;
  \texttt{gate = false} opens it at once.
\item
  Side-lane nodes are always eligible and never block or unlock anything
  on the main path.
\item
  \texttt{hack} with one main and one side candidate prints both (Next +
  Quick win); \texttt{hack 10} prefers whichever fits.
\item
  Active days per level counts distinct logical days, not entries (two
  entries on one day = 1).
\item
  \texttt{lint} rejects a node whose \texttt{level} does not exist and a
  pack with no main level.
\item
  No engine source file contains the strings "bandit", "juice",
  "overthewire" (grep test).
\end{itemize}

\hypertarget{hackstreak-spec-v2-13-prompt-for-the-coding-agent}{%
\subsubsection{13. Prompt for the coding
agent}\label{hackstreak-spec-v2-13-prompt-for-the-coding-agent}}

\begin{quote}
Read \path|hackstreak-spec-v2.md| (it supersedes v1 where they
conflict; also read v1 for the streak rules and algorithm). Build
\textbf{Phase 1 AND Phase 1b of v2}: Python 3.11, standard library only
for Phase 1 (Phase 1b\textquotesingle s Telegram call may use
\texttt{urllib}, still no third-party dependency), one installable CLI
\texttt{hack}. Write the tests from v1 section 11 and v2 section 12
FIRST (pytest), then implement until they pass. Install
\path|packs/home-hack-lab.toml| as the sample pack and run
\texttt{hack pack lint} on it. Ask me before deviating from the spec.
When finished show: file tree, \texttt{pytest} output, and a demo of
\texttt{hack pack install}, \texttt{hack mission start},
\texttt{hack backfill}, \texttt{hack levels}, \texttt{hack},
\texttt{hack 20}, \path|hack done ... ---complete|,
\texttt{hack stats}, \path|hack config nudge-test| (or an
explanation of what is stubbed if Telegram credentials are not available
in this environment).
\end{quote}
